\documentclass[10pt,a4paper]{article}
\usepackage[utf8]{inputenc}
\usepackage{amsmath,amssymb,amsfonts}
\usepackage{graphicx}
\usepackage{hyperref}

\title{Urban Resilience Through Cognitive Computing Systems}
\author{Kenan Degirmenci \and Kevin C. Desouza \and Richard T. Watson}
\date{\today}

\begin{document}
\maketitle

Journal of Urban Technology

ISSN: 1063-0732 (Print) 1466-1853 (Online) Journal homepage: www.tandfonline.com/journals/cjut20

Urban Resilience Through Cognitive Computing Systems

Kenan Degirmenci, Kevin C. Desouza \& Richard T. Watson

To cite this article: Kenan Degirmenci, Kevin C. Desouza \& Richard T. Watson (2024) Urban Resilience Through Cognitive Computing Systems, Journal of Urban Technology, 31:2, 117-124, DOI: 10.1080/10630732.2024.2317767

To link to this article: https://doi.org/10.1080/10630732.2024.2317767

© 2024 The Author(s). Published by Informa UK Limited, trading as Taylor \& Francis Group

Published online: 27 Mar 2024.

Submit your article to this journal

Article views: 1235

View related articles

View Crossmark data

Citing articles: 1 View citing articles

Full Terms \& Conditions of access and use can be found at https://www.tandfonline.com/action/journalInformation?journalCode=cjut20


\section{COMMENTARY}

Urban Resilience Through Cognitive Computing Systems

Kenan Degirmenci a, Kevin C. Desouza b, and Richard T. Watson c

aQueensland University of Technology, Faculty of Science, School of Information Systems, 2 George Street, Brisbane QLD 4000, Australia; bQueensland University of Technology, Faculty of Business and Law, School of Management, 2 George Street, Brisbane QLD 4000, Australia; cDigital Frontier Partners, 4/350 Collins Street, Melbourne VIC 3000, Australia

ABSTRACT Global urbanization and heat-related fatalities are rapidly increasing, while the full potential of advanced information technologies to mitigate urban heat has not yet been used by city planners and policymakers. Cognitive computing systems (CCSs), which mimic human information processing and reasoning abilities, can transform how cities strategize and operate. In this commentary, we develop a framework and outline future research directions to understand the disconnect between a strategic and operational CCS, which can facilitate an effective interplay between urban policy (e.g., framing alternative policies) and urban technology (e.g., engaging local communities with data generated by sensor networks).

KEYWORDS urban resilience; urban policy; urban technology; cognitive computing systems (CCS); sustainability

Urban Heat as a Major Challenge to Making Cities More Resilient

Science tells us that to keep global temperatures from rising by more than 1.5°C, we need to achieve net-zero emissions by mid-century. Sustainable and equitable urban cooling must be a part of cities’ efforts to reach net-zero energy targets.—Inger Andersen, Executive Direc- tor of the UN Environment Programme (2021)

Rapid global urbanization has increased the vulnerability of cities, which now face more frequent heatwaves, storms, and floods (Selby and Desouza, 2019). The July–August 2003 European heatwave caused 72,210 fatalities from heat-related deaths in major European cities; the economic damage accounted for over US\$12.1 billion (EM-DAT, 2024). More recently, extreme heatwaves have caused wildfires, droughts, and deaths globally, includ- ing in Europe, Africa, Asia, and America (NASA, 2022). For the first time in 2022, the UK Meteorological Office issued a red warning for heat, its most extreme alert (Cappucci and Samenow, 2022).

Heat kills more people than any other extreme weather event, because heatwaves are now more intense (Calkins et al., 2016). Urban heat islands (UHI) in metropolitan areas

© 2024 The Author(s). Published by Informa UK Limited, trading as Taylor \& Francis Group This is an Open Access article distributed under the terms of the Creative Commons Attribution-NonCommercial-NoDerivatives License (http://creativecommons.org/licenses/by-nc-nd/4.0/), which permits non-commercial re-use, distribution, and reproduction in any medium, provided the original work is properly cited, and is not altered, transformed, or built upon in any way. The terms on which this article has been published allow the posting of the Accepted Manuscript in a repository by the author(s) or with their consent.

CONTACT Kevin C. Desouza kevin.c.desouza@gmail.com Queensland University of Technology, Faculty of Business and Law, School of Management, 2 George Street, Brisbane QLD 4000, Australia.

JOURNAL OF URBAN TECHNOLOGY 2024, VOL. 31, NO. 2, 117–124 https://doi.org/10.1080/10630732.2024.2317767

can cause city temperatures to be 2–12°Chigher than surrounding rural areas (Li et al., 2019). More than 40 million people in the United States live in UHI regions, significantly affecting health and equity (Sharma, 2023). People residing in low-income neighbor- hoods are particularly stuck in UHI regions compared to their wealthier counterparts (Newsome, 2023). The UHI effect is strongest in densely built-up areas, which absorb heat because vegetation and water bodies have been replaced by impervious materials, such as pavements, concrete buildings, and asphalt roads (Estrada et al., 2017). For cities in hot, warm, or temperate climates, UHI can increase energy consumption, elevate emissions of air pollutants and greenhouse gases, compromise human health and comfort, and impair water quality (EPA, 2023).1

Natural hazards such as floods and heatwaves often result in substantial immediate destruction of resources and the disruption of services (Desouza and Flanery, 2013). However, climate deterioration is a function of slow-burning factors that seldom get the requisite attention because their daily impact is often indiscernible. Advanced infor- mation technologies, which play an increasingly transformative role in city operations (Chin and Guthrie, 2023), can alleviate intergenerational discounting, where investments for urban cooling and climate improvement will not generate benefits for another gen- eration (Jacquet et al., 2013). Society must react to exogenous shocks and make the necessary changes to address threats to resilience proactively. Advanced information systems can play a vital role in enabling urban planners to assess, analyze, simulate, and visualize the impacts of their plans on UHI.

Attempts to address UHI reveal the effects of conflicting priorities between social, environmental, and economic interests. For example, low-income neighborhoods often have higher building densities and fewer vegetative resources, such as parks. This exposes their residents to higher health hazards during periods of extreme heat (Hsu et al., 2021) and energy demand for cooling is usually higher in these areas. Urban resili- ence requires policies that mitigate UHI by balancing the demand for buildings with more green spaces (Tuczek et al., 2022).2 However, the ramifications of such policies can take years or even decades to unfold due to the long-term nature of climate impacts.

Urban planning policies to maximize the desired impacts and minimize the undesir- able outcomes can exploit technological advances such as artificial intelligence, predictive modelling, and environmental monitoring (Corbett and Mellouli, 2017). Smart cities are already laden with various technologies to collect and analyze unprecedented data volumes (Brandt et al., 2018). However, current practices fail to harness their full poten- tial. One general class of technologies, cognitive computing systems (CCSs),3 which can mimic human information processing and reasoning capabilities, can transform how cities strategize and operate.

Integrating Cognitive Computing Systems in Urban Design Problems

Growing urbanization increases the complexity of cities, which require advanced infor- mation systems to achieve efficient real-time matching and coordination of tasks and resources to enable effective demand and supply responses (Ketter et al., 2020). CCSs can address such needs because they can mimic five fundamental systems typically deployed by organizations to create value (Watson and Pitt, 2022): systems of framing and inquiry for strategic purposes, systems of engagement and production for

118 K. DEGIRMENCI ET AL.

operational procedures, and systems of record at the intersection to enable strategic and operational coordination (See Table 1).

Systems of engagement support collaboration and coordination to achieve shared goals, while systems of production allow humans to create products and services. Systems of framing are used to justify human behavior and decision-making, and systems of inquiry generate knowledge using replicable methods. Systems of record allow storage and retrieval of data (Watson, 2021). While humans have developed these five systems over many years, CCSs now incorporate recent digital technologies, such as cloud comput- ing, the Internet of Things (IoT), machine learning, and big data analytics to save time and effort and reduce human error (Gupta et al., 2018). Early CCSs include IBM’s Deep Blue, chess-playing software that defeated world champion Garry Kasparov in 1997 (although it lost the first match in 1996), and Apple’s Siri, a virtual assistant introduced in 2011 that uses natural language processing to process voice commands. Today, CCSs are more advanced, as we see with chatbots based on large language models.

A CCS can support an effective interplay between urban policy and technology (See Figure 1), alleviating urban heat and leading to desirable social, environmental, and econ- omic outcomes.

The Digital Urban Climate Twin (DUCT), currently being developed as part of the Cooling Singapore project (Ruefenacht and Acero, 2017), is an exemplar CCS. This digital twin allows urban planners to conduct what-if analyses (system of inquiry) to evaluate alternative policies for heat mitigation strategies and justify their decisions (system of framing).4 Another example is the National Integrated Heat Health Infor- mation System (NIHHIS).5 Acollaboration with Climate Adaptation Planning and Ana- lytics (CAPA) Strategies has yielded Heat Watch, a citizen science project. Residents are equipped with temperature and humidity sensors to attach to their cars and bikes (system of production) for generating fine-granular heat mapping (system of record). Heat maps, created in over 50 US cities, help to identify needs for extreme heat services (system of production). DUCT and NIHHIS require sensor networks to collect and process data for informed decision-making (system of inquiry).

There is often a disconnect between a strategic and operational CCS. However, there is potential to facilitate the interplay between urban policy (e.g., framing alternative policies using digital twins) and urban technology (e.g., engaging local communities to create heat maps). The conjunctive use of urban policy and technology (Degirmenci et al., 2021) afforded by a joint strategic and operational CCS can mitigate urban heat, and, in turn, alleviate social vulnerability, environmental burdens, and economic loss (See Figure 1). There are three areas in our framework: (1) a CCS facilitating the interplay

Table 1. Types of systems

Type Purpose Mechanisms

Engagement Collaborating and coordinating activities Touch, gesturing, speech, and writing Production Creating products and services Standard operating procedures to increase

efficiency and product consistency Framing Justifying the reason for behaving in a particular

way or promoting a particular opinion

The use of mental models to focus discussion, justify

a viewpoint, and generalize knowledge Inquiry Generating knowledge Use of replicable methods to justify or validate an

inference Record Recording and retrieving data Tables, ledgers, and written reports Source: Watson, 2021

JOURNAL OF URBAN TECHNOLOGY 119

between urban policy and technology; (2) the conjunctive use of urban policy and tech- nology to mitigate urban heat; and (3) outcomes of urban heat as a triple-bottom-line responsibility: social (people), environment (planet), and economic (profit). When com- puting economic outcomes, it is crucial to recognize that public losses (externalities such as carbon emissions) are partially or fully paid for by public financial means, which are funded by taxes (a private loss) on personal gains.6 There can also be a gain from improv- ing the ecological desirability of an area because of cooler temperatures, cleaner air, and environmental attractiveness. The goal is to provide a societal gain to get broad buy-in and lessen the power of vested interests to prevent or delay UHI mitigation.

Future Directions of Cognitive Computing Research for UHI Mitigation

There are various ways to link CCSs directly to the three outcomes. For example, when it comes to engagement, CCSs can focus on technologies like edge computing, sensors, and (social) platforms to build connections across layers of society, from individuals to organizations, on UHI issues. With framing, a CCS concentrates on simulating the con- sequences of design choices and their outcomes in the near and long term because decision-makers, the framers of action, understand the costs and benefits. For pro- duction, a CCS can focus on technologies like robotic process automation to streamline the environmental protocols for new developments. Inquiry involves methods like deep learning to discover unknown patterns, which can be linked to digital twins and heat maps to investigate urban fragilities in real-time. Finally, for records, sensor networks are required to collect and process data to feed a CCS. To guide future studies on urban resilience through CCSs, we present 15 research questions (See Table 2) covering

Figure 1. Framework to mitigate UHI through cognitive computing systems

120 K. DEGIRMENCI ET AL.

Table 2. Future research guide

Social Environment Economic Study suggestion Key approach

Engagement RQ1: What types of CCSs

effectively engage

community leaders in

reducing UHI?

RQ2: How can a CCS

nudge residents toward

more sustainable

behaviors?

RQ3: How can a CCS inform investors

about the losses due to UHI? Implement system dynamics and

agent-based simulations to

analyze the overall behavior of

urban systems and interactions of

individual agents

Establish intelligent

communication systems that

disseminate critical information

effectively

Production RQ4: How do we introduce

CCS-driven procedures to

alert urban planners to

social vulnerabilities?

RQ5: How do we

introduce CCS-driven

procedures to address

urban design problems?

RQ6: How do we introduce CCS-driven

procedures to alert urban planners

to economic vulnerabilities?

Perform action design research to

improve knowledge

dissemination through CCS-

driven procedures for urban

planning

Create digital solutions by

integrating a CCS into a cohesive

platform fostering collaboration

among stakeholders

Framing RQ7: How can a CCS be

used to design policies to

justify distributive energy

justice?

RQ8: What form of a CCS

will help residents

identify practical

actions for reducing


\section{UHI?}

RQ9: How do we create a CCS that

enables key stakeholders to

understand sustainable urban

planning possibilities and reduce

intergenerational discounting?

Conduct experimental studies to

evaluate the effectiveness of CCS

implementation for urban policy

and technology interventions

Develop policies to enhance CCS

implementation into existing

and new urban resilience

frameworks

Inquiry RQ10: How can we design a

CCS to identify social

vulnerabilities?

RQ11: How can UHI

modelling be improved

with a CCS?

RQ12: What type of CCS would nudge

urban planners and investors to

lower UHI outcomes?

Undertake data mining studies for

predictive modelling based on

historical and real-time data

Identify vulnerable areas and

analyze UHI trends and

consumption patterns

Record RQ13: What demographic

and psychographic data

are needed to build a


\section{CCS?}

RQ14: What data does a

CCS require to estimate

UHI environmental

outcomes?

RQ15: What economic and natural

capital data are needed to build a


\section{CCS?}

Execute big data analytics to

provide a structured and

organized representation of

relevant data for CCS integration

Apply database management,

systems design and

implementation, and data

analytics for integrated CCS

functionality and performance

JOURNAL OF URBAN TECHNOLOGY 121

the three triple-bottom-line impact factors and the five CCS areas, study suggestions, and key approach considerations.

Urban studies disciplines can contribute significantly by creating evidence-driven sol- utions to tackle UHI. Doing so will require collaboration with scholars from domains such as information systems, public policy, economics, and engineering. Citizen science and community engagement are areas that can offer existing frameworks and theories to build solutions around a CCS for engagement. Science communication and data visualization are natural allies to inform a CCS for framing. The urban discipline could co-create a CCS for production and inquiry by working with environmental scientists and sustainability advo- cates. Given the significant expertise within the urban discipline around smart cities, IoT in urban environments, and geographical information systems (GIS), it should lead the con- versation for creating standards to promote the interoperability of a CCS focused on records.

Notes

1. In cold climates, UHI can have positive effects, such as reduced demand for domestic heating (Roth, 2013). 2. The study proposed an optimization model for Brisbane in Australia and showed that an increase of buildings and decrease of vegetation coverage led to revenue growth by about three billion US dollars based on land valuation data, but the maximum UHI intensity level rose from 4°Cto 5°C (Tuczek et al., 2022). 3. Cognitive computing is “a discipline that links together neurobiology, cognitive psychology, and artificial intelligence” (Valiant, 1995: 2). 4. Strategies for mitigating urban heat include increasing vegetation (e.g., trees) and water body coverage (e.g., ponds on roofs); optimizing urban geometry (e.g., varying building forms to improve wind capture); using light-colored surfaces and reflective materials; increasing shading, blinds, and shutters; boosting public transport use and promoting active mobility like cycling; lowering energy consumption by using energy-efficient equip- ment (Roth, 2013; Ruefenacht and Acero, 2017). 5. https://www.heat.gov/ 6. The costs of many externalities, such as carbon emissions, will be paid by future generations.

Notes on Contributors

Kenan Degirmenci is a lecturer in the School of Information Systems at Queensland University of Technology. His research focuses on technology adoption, particularly within the energy sector.

Kevin C. Desouza is a Professor of Business, Technology, and Strategy in the School of Manage- ment at Queensland University of Technology. He is a Non-resident Senior Fellow in the Govern- ance Studies Program at the Brookings Institution.

Richard T. Watson is the research director for Digital Frontier Partners and also Regents Professor and the J. Rex Fuqua Distinguished Chair for Internet Strategy Emeritus of the University of Georgia. He is a Schüller Senior Scholar at the University of Erlangen-Nuremberg and an honorary visiting professor at the Queensland University of Technology. He is a former President of the Association for Information Systems and was awarded its highest honor, a LEO, for his achieve- ments in information systems.

Disclosure Statement

No potential conflict of interest was reported by the author(s).

122 K. DEGIRMENCI ET AL.


\section{ORCID}

Kenan Degirmenci http://orcid.org/0000-0002-4046-4526 Kevin C. Desouza http://orcid.org/0000-0002-4734-3081 Richard T. Watson http://orcid.org/0000-0003-0664-8337

References

T. Brandt, W. Ketter, L. M. Kolbe, D. Neumann, and R. T. Watson, “Smart Cities and Digitized

Urban Management,” Business \& Information Systems Engineering 60: 3 (2018) 193–195. M. M. Calkins, T. B. Isaksen, B. A. Stubbs, M. G. Yost, and R. A. Fenske, “Impacts of Extreme Heat

on Emergency Medical Service Calls in King County, Washington, 2007–2012: Relative Risk and Time Series Analyses of Basic and Advanced Life Support,” Environmental Health 15 (2016) 1–13. M. Cappucci and J. Samenow, “These Maps Show How Excessively Hot it is in Europe and the U.S.,”

The Washington Post (July 18, 2022), Accessed January 8, 2024. J. T. Chin and A. Guthrie, “What Makes a City ‘Smart’ and Who Decides? From Vision to Reality

in the USDOT Smart City Challenge,” Journal of Urban Technology 30: 4 (2023) 3–31. J. Corbett and S. Mellouli, “Winning the SDG Battle in Cities: How an Integrated Information

Ecosystem Can Contribute to the Achievement of the 2030 Sustainable Development Goals,” Information Systems Journal 27: 4 (2017) 427–461. K. Degirmenci, K. C. Desouza, W. Fieuw, R. T. Watson, and T. Yigitcanlar, “Understanding Policy

and Technology Responses in Mitigating Urban Heat Islands: A Literature Review and Directions for Future Research,” Sustainable Cities and Society 70 (2021) 1–35. K. C. Desouza and T. H. Flanery, “Designing, Planning, and Managing Resilient Cities: A

Conceptual Framework,” Cities 35 (2013) 89–99. EM-DAT, “The International Disaster Database,” Centre for Research on the Epidemiology of

Disasters (January 4, 2024), Accessed January 8, 2024. EPA, “Heat Island Impacts,” United States Environmental Protection Agency (August 28, 2023),

Accessed January 8, 2024. F. Estrada, W. J. W. Botzen, and R. S. J. Tol, “A Global Economic Assessment of City Policies to

Reduce Climate Change Impacts,” Nature Climate Change 7: 6 (2017) 403–406. S. Gupta, A. K. Kar, A. Baabdullah, and W. A. Al-Khowaiter, “Big Data with Cognitive Computing:

A Review for the Future,” International Journal of Information Management 42 (2018) 78–89. A. Hsu, G. Sheriff, T. Chakraborty, and D. Manya, “Disproportionate Exposure to Urban Heat

Island Intensity across Major US Cities,” Nature Communications 12 (2021) 1–11. J. Jacquet, K. Hagel, C. Hauert, J. Marotzke, T. Röhl, and M. Milinski, “Intra- and Intergenerational

Discounting in the Climate Game,” Nature Climate Change 3 (2013) 1025–1028. W. Ketter, B. Padmanabhan, G. Pant, and T. S. Raghu, “Special Issue Editorial: Addressing Societal

Challenges through Analytics: An ESG ICE Framework and Research Agenda,” Journal of the Association for Information Systems 21: 5 (2020) 1115–1127. D. Li, W. Liao, A. J. Rigden, X. Liu, D. Wang, S. Malyshev, and E. Shevliakova, “Urban Heat Island:

Aerodynamics or Imperviousness?,” Science Advances 5: 4 (2019) 1–4. NASA, “Heatwaves and Fires Scorch Europe, Africa, and Asia,” Earth Observatory (July 13, 2022),

Accessed January 8, 2024. M. Newsome, “Discrimination Has Trapped People of Color in Unhealthy Urban ‘Heat Islands’,”

Scientific American (September 19, 2023), Accessed January 8, 2024. M. Roth, “Urban Heat Islands,” in H. J. Fernando, ed., Handbook of Environmental Fluid

Dynamics: Systems, Pollution, Modeling, and Measurements (Boca Raton: CRC Press, 2013) 143–159.

JOURNAL OF URBAN TECHNOLOGY 123

L. A. Ruefenacht and J. A. Acero, “Strategies for Cooling Singapore,” Singapore-ETH Centre (June

1, 2017), Accessed January 8, 2024. J. D. Selby and K. C. Desouza, “Fragile Cities in the Developed World: A Conceptual Framework,”

Cities 91 (2019) 180–192. N. Sharma, “More Than 40 Million People in the U.S. Live in Urban Heat Islands, Climate Group

Finds,” NBC News (July 26, 2023), Accessed January 8, 2024. M. Tuczek, K. Degirmenci, K. C. Desouza, R. T. Watson, T. Yigitcanlar, and M. H. Breitner,

“Mitigating Urban Heat with Optimal Distribution of Vegetation and Buildings,” Urban Climate 44 (2022) 1–8. UN Environment Programme, “UN Issues New Guidance to Address Warming in Cities,” United

Nations (November 3, 2021), Accessed January 8, 2024. L.G. Valiant, “Cognitive computation,” Proceedings of IEEE 36th Annual Foundations of Computer

Science, Milwaukee, WI, USA, 1995, pp.: 2–3, doi: 10.1109/SFCS.1995.492456. R. T. Watson, Capital, Systems, and Objects: The Foundation and Future of Organizations

(Singapore: Springer, 2021). R. T. Watson and L. F. Pitt, “Transcendent Service Management,” Journal of Service Management

33: 1 (2022) 1–8.

124 K. DEGIRMENCI ET AL.

\end{document}